% secV4_coverage.tex -- Section 6. CONTAINER (27 Oct 2026): text to be written by the author.
% Word budget: 400 words (JPART count: running text, footnotes and citations as printed).
% Outline, numbers and sources: referee_EE_2026-10/ESQUELETO.md, section 6. The archived text
% (_arquivo_texto_IA_2026-10-27/) is not a source.

\section{Coverage of the request record}\label{sec:coverage}

The request record underpins the cohort trends analysed below and the proposed measure of application flows, but not the Goal~5 compliance exercise: every specification in our measurement grid relies on published aggregate waiting-list counts rather than individual application records. The supplied files are not SME's live system. Stocks reconstructed from them for 2010 reach only 20--33\% of the published list (Table~\ref{tab:recon}), which indicates incomplete coverage in the supplied records, though not necessarily in the underlying administrative system.

We assess coverage independently by matching first requests to childcare entrants in Escola On-Line (EOL) enrolment records. The matched share is 42.1\%, 47.9\%, and 57.2\% in 2010, 2011, and 2012, respectively, against 98.6--99.9\% in 2013--2018 (Table~\ref{tab:coverage}; Figure~\ref{fig:coverage}). The discontinuity is concentrated in February: only 41.4\% of 87,228 February 2012 entrants have a matching request, versus 98.3\% of 33,773 entrants in other months. Across centre types, February coverage rises from 32.0--43.6\% in 2012 to 95.8--98.8\% in 2013 (Table~\ref{tab:covtype}), and, among the 1,339 units recording 20 or more entrants in each February, 92.4\% improve by at least 30 percentage points. This simultaneous change is consistent with a boundary in the supplied files; its administrative cause is not identified.

We define the window of complete coverage with a prespecified rule: completeness and referral recording may each fall at most two percentage points below their 2014--2016 averages, yielding thresholds of 96.4\% and 97.4\% (Table~\ref{tab:completeness}). Completeness by semester of first request rises from 92.8\% in the first half of 2012 to 98.0\% in the second and stays between 97.9\% and 99.9\% through 2017. We therefore restrict analyses of first requests to July 1, 2012 onward. The records are predominantly prospective: among 2013 entrants, the median request precedes enrolment by 324 days, 89.5\% precede it by more than 30 days, and only 0.64\% follow it (Table~\ref{tab:prospect}).

\subsection{Consequences for cohort trends}\label{sec:naive}

Ignoring this window distorts cohort comparisons (Figure~\ref{fig:level}). For children born in 2010, enrolment before 24 months is 7.4\% of births in the request record (Table~\ref{tab:enrolledbyage}) but 25.1\% in EOL enrolment records. The two sources converge for the 2012 cohort (33.2\% versus 31.9\%) and differ by only 0.5--1.0 percentage points in 2013--2016. Without imputation, 39\% of the apparent 2010--2016 increase in the record-based measure reflects changes in record coverage, not in behaviour (Table~\ref{tab:level}). This limitation concerns analyses based on the request record, not the Goal~5 classifications.
